\documentclass[11pt,a4paper]{article}
\usepackage[utf8]{inputenc}
\usepackage[T1]{fontenc}
\usepackage[spanish,es-noshorthands]{babel}
\usepackage{lmodern}
\usepackage{microtype}
\usepackage[margin=2.3cm]{geometry}
\usepackage{amsmath,amssymb}
\usepackage{booktabs,longtable,array,tabularx}
\usepackage{enumitem}
\usepackage[dvipsnames]{xcolor}
\usepackage{tikz}
\usetikzlibrary{arrows.meta,positioning,shapes.geometric,fit,calc}
\usepackage[colorlinks=true,linkcolor=NavyBlue,urlcolor=NavyBlue]{hyperref}
\usepackage{titlesec}
\usepackage{fancyhdr}

\definecolor{accent}{HTML}{1F4E79}
\definecolor{boxbg}{HTML}{EEF3F8}
\titleformat{\section}{\Large\bfseries\color{accent}}{\thesection}{0.6em}{}
\titleformat{\subsection}{\large\bfseries\color{accent}}{\thesubsection}{0.6em}{}
\pagestyle{fancy}\fancyhf{}
\lhead{\small PISR-NSGA-II --- documento técnico}\rhead{\small\thepage}
\setlength{\parskip}{4pt}\setlength{\parindent}{0pt}
\setlist{itemsep=2pt,topsep=3pt}

\newcommand{\code}[1]{\texttt{\small #1}}
\newcommand{\nota}[1]{\par\smallskip\noindent\colorbox{boxbg}{\parbox{\dimexpr\linewidth-2\fboxsep}{\small #1}}\par\smallskip}

\title{\color{accent}\bfseries PISR-NSGA-II\\[4pt]\large Documento técnico del procedimiento completo}
\author{Referencia personal del autor --- rama \code{grade}}
\date{Noviembre 2026}

\begin{document}
\sloppy
\maketitle
\thispagestyle{fancy}

\begin{abstract}
Este documento describe, paso a paso, qué hace PISR-NSGA-II desde que recibe una ecuación diferencial
hasta que entrega una fórmula: las verificaciones y análisis que corren \emph{antes} de la búsqueda, la
representación de las soluciones, cómo se evalúa cada individuo, qué ocurre en cada generación, cómo se
elige y pule el campeón, y qué se hace con los resultados después. Cada afirmación remite al archivo del
código donde está implementada. Los valores de los parámetros son los de \code{src/pi\_solver.cpp}
(espacio de nombres \code{Config}).
\end{abstract}

\tableofcontents
\newpage

%=====================================================================
\section{Visión general}
%=====================================================================

\paragraph{Problema.} Dada una ecuación diferencial $\mathcal{N}[u] = 0$ en $\Omega=[0,1]^d$ (y
$t\in[0,1]$ si no es estacionaria) con sus datos de borde o iniciales, encontrar una \emph{fórmula
cerrada} $u(x,y,t)$ que la satisfaga. El algoritmo usa \textbf{solo} el residuo de la ecuación y los
datos de borde; la solución de referencia nunca entra en la búsqueda y se usa únicamente para medir el
resultado al final.

\paragraph{Tesis.} Un físico entiende mejor una fórmula corta que una grilla de números. PISR-NSGA-II
cambia algo de precisión por interpretabilidad: busca la expresión más simple que resuelva la ecuación,
y si existe una forma cerrada, la recupera.

\paragraph{Idea del método.} Programación genética (árboles de expresión) con NSGA-II sobre tres
objetivos a minimizar:
\begin{align*}
f_1 &= \text{residuo de la EDP en el dominio},\\
f_2 &= \text{error de los datos de borde},\\
f_3 &= \text{tamaño del árbol (nodos)}.
\end{align*}
Las derivadas se calculan exactamente con diferenciación automática hacia adelante (jets) sobre
aritmética compleja.

\begin{figure}[h]
\centering
\begin{tikzpicture}[node distance=5mm and 4mm, every node/.style={font=\small},
  box/.style={draw=accent, rounded corners=2pt, fill=boxbg, align=center, minimum height=8mm, text width=24mm},
  arr/.style={-{Stealth[length=2mm]}, accent, thick}]
\node[box] (ref) {Referencia\\numérica\\(solo evaluación)};
\node[box, right=of ref] (orc) {Oráculo\\de coherencia};
\node[box, right=of orc] (pri) {Sondeo del\\operador\\(priors)};
\node[box, right=of pri] (ini) {Población\\inicial sesgada};
\node[box, below=9mm of ini] (loop) {Bucle evolutivo\\NSGA-II\\(hasta 1000 gen.)};
\node[box, left=of loop] (hof) {Campeón por\\validación física\\(Hall of Fame)};
\node[box, left=of hof] (pol) {Pulido final\\de constantes};
\node[box, left=of pol] (out) {Fórmula, grilla,\\frentes, registro};
\draw[arr] (ref) -- (orc); \draw[arr] (orc) -- (pri); \draw[arr] (pri) -- (ini);
\draw[arr] (ini) -- (loop); \draw[arr] (loop) -- (hof); \draw[arr] (hof) -- (pol); \draw[arr] (pol) -- (out);
\end{tikzpicture}
\caption{Flujo por problema y semilla (\code{run\_once} en \code{main.cpp}).}
\end{figure}

\paragraph{Ejecución.} \code{./run\_tests.sh} compila y corre las pruebas (gradiente y oráculo de los 21
problemas); luego \code{./build\_release/PISR-NSGA-II -{}-cores N -{}-replicable -{}-runs 30} corre la suite.
Cada semilla escribe en \code{results/run\_<k>/}; \code{-{}-only <EDP>} corre un solo problema y
\code{-{}-no-hybrid} la ablación sin representación híbrida (en \code{results/ablation\_no\_hybrid/}).

%=====================================================================
\section{La suite de problemas}
%=====================================================================
Definida en \code{benchmark\_suite()} (\code{src/pde\_problems.cpp}) y replicada para el reporte en
\code{report/benchmarks.py}. Residuos en \code{PDEProblem::pde\_residual\_ad}.

{\small
\begin{longtable}{@{}p{3.1cm}p{5.0cm}p{2.9cm}p{3.3cm}@{}}
\toprule
\textbf{Problema} & \textbf{Ecuación} & \textbf{Datos} & \textbf{Referencia}\\\midrule\endhead
Laplace 1D & $u''=0$ & $u(0),u(1)$ & $x$\\
Poisson 1D & $u''=-\pi^2\sin\pi x$ & $u(0),u(1)$ & $\sin\pi x$\\
Oscilador armónico 1D & $u''-x^2u+u=0$ & $u(0),u(1)$ & $e^{-x^2/2}$\\
Airy 1D / 2D & $u''=xu$;\ $\nabla^2u=(x+y)u$ & borde & $\mathrm{Ai}(x)$;\ $\mathrm{Ai}(2^{-1/3}(x+y))$\\
Fisher--KPP 1D / 2D & $\nabla^2u+u(1-u)=0$ & borde & numérica\\
Duffing 1D / 2D & $\nabla^2u+u+u^3=\sin\pi x\,(\sin\pi y)$ & borde & numérica\\
Thomas--Fermi 1D / 2D & $\nabla^2u=u^{3/2}/\sqrt r$ & $u(0)=1,\ u(1)=0.424$ & numérica\\
Schrödinger 1D / 2D & $\nabla^2u+4d\,u=0$ (complejo) & borde & $e^{2ix}$;\ $e^{2i(x+y)}$\\
Schrödinger no estac. 1D / 2D & $i u_t+\nabla^2u=0$ & $u(\cdot,0)$ y borde $\forall t$ & $e^{i\pi(x-\pi t)}$;\ $e^{i\pi(x+y-2\pi t)}$\\
Navier--Stokes (Kovasznay) & $\psi_y\nabla^2\psi_x-\psi_x\nabla^2\psi_y=\nu\nabla^4\psi$, $\nu=0.05$ & $\psi,\ \partial_n\psi$ & $y-\tfrac{1}{2\pi}e^{\lambda x}\sin 2\pi y$\\
Navier--Stokes (Taylor--Green) & $\omega_t+\mathbf u\cdot\nabla\omega=\nu\nabla^2\omega$, $\nu=0.01$ & $\psi(\cdot,0)$, $\psi,\partial_n\psi$ & $\sin\pi x\sin\pi y\,e^{-2\pi^2\nu t}$\\
Lane--Emden ($n=3$) & $u''+\tfrac2x u'+u^3=0$ & IVP: $u(0)=1,u'(0)=0$ & numérica\\
Troesch ($\lambda=3$) & $u''=3\sinh 3u$ & $u(0)=0,u(1)=1$ & numérica\\
Ginzburg--Landau & $u''+u-u^3=0$ & $u(0),u(1)$ & $\tanh(x/\sqrt2)$\\
Painlevé I & $u''=u^2+x$ & $u(0)=0,u(1)=\tfrac12$ & numérica\\\bottomrule
\end{longtable}}

Decisiones de planteamiento que conviene recordar:
\begin{itemize}
\item \textbf{Schrödinger} usa $k^2=4$ (no $\pi^2$) para evitar la resonancia de Dirichlet: con $\pi^2$
  el problema tenía infinitas soluciones.
\item \textbf{Lane--Emden} es un problema de valor inicial ($u(0)=1$, $u'(0)=0$), no de contorno.
\item \textbf{Thomas--Fermi} está truncado a $[0,1]$ con $u(1)=0.424$, el valor tabulado de la solución
  semi-infinita; ese dato es parte del planteamiento, no una fuga.
\item \textbf{Ginzburg--Landau}: la solución exacta es $\tanh(x/\sqrt2)$ (antes se usaba $\tanh x$, que
  no resuelve la ecuación).
\item \textbf{Navier--Stokes} usa la función de corriente y el residuo real de orden 4 (antes había
  proxies de orden 2). Kovasznay con $\mathrm{Re}=20$; por ser de orden 4, la frontera incluye
  $\partial_n\psi$.
\item Los problemas \textbf{no estacionarios} dan la frontera espacial para todo $t$ y la condición
  inicial en $t=0$; sin la condición inicial quedarían indeterminados.
\end{itemize}

%=====================================================================
\section{Antes de la búsqueda}
%=====================================================================
Todo lo de esta sección ocurre por problema, antes de crear la población (\code{run\_once} en
\code{main.cpp}, y el constructor \code{PISolver::PISolver}).

\subsection{Referencia numérica (solo para evaluar)}
Para los problemas sin forma cerrada, \code{NumericalSolver::reference} (\code{src/numerical\_solver.cpp})
calcula una solución de alta precisión usando la \emph{misma} ecuación (\code{pde\_residual\_ad}) y los
mismos datos de borde que recibe el algoritmo:
\begin{itemize}
\item \textbf{1D de contorno:} Newton amortiguado sobre diferencias finitas centradas de orden 2;
  jacobiano tridiagonal por diferencias finitas, resuelto con el algoritmo de Thomas.
\item \textbf{1D de valor inicial} (Lane--Emden): RK4 arrancando en $x_0>0$ con una serie de Taylor
  consistente ($u''(0)$ se resuelve por secante sin conocer la solución), porque el origen es singular.
\item \textbf{2D de Dirichlet:} SOR no lineal (Newton por punto), arrancando de la interpolación de
  Coons de los datos.
\end{itemize}
El reporte del solver incluye el residuo discreto y la diferencia de Richardson entre resoluciones $h$
y $h/2$ (estimación del error de discretización). Si no converge, el problema no se corre.

\nota{La referencia solo la leen: el error del campeón al final, el registro de convergencia (columna
\code{champion\_ref\_err}, solo escritura) y el oráculo. Ninguna decisión de la búsqueda la usa.}

\subsection{Oráculo de coherencia}
\code{run\_oracle} (\code{src/oracle.cpp}) verifica que el problema esté bien planteado y que la
maquinaria funcione. Si falla, el problema se omite (salvo \code{-{}-skip-oracle}). Comprobaciones:
\begin{enumerate}
\item \textbf{La referencia satisface la EDP}, con un jet por diferencias finitas. Pasa si el residuo
  RMS es menor que la tolerancia o, para referencias analíticas, si duplicar el paso multiplica el
  residuo por entre 3 y 5 (es solo error de truncación $O(h^2)$, no una referencia equivocada).
\item \textbf{La referencia satisface los datos de borde e iniciales.}
\item \textbf{Bien planteado:} el solver numérico, solo con la ecuación y el borde, reproduce la solución
  exacta.
\item \textbf{El árbol de la solución conocida es factible y casi exacto} dentro del algoritmo, en las
  tres representaciones: sin ansatz, híbrida (U directo) y con ansatz ($N=(u-L)/B$). Esto detecta
  restricciones de factibilidad que excluirían la respuesta correcta.
\item \textbf{Priors:} el orden de derivación detectado coincide con el declarado, y se informa qué
  propiedades del operador se cumplen también en la solución (diagnóstico).
\end{enumerate}
\code{tests/oracle\_test} corre el oráculo de los 21 problemas (21/21 aprobados) y
\code{tests/gradient\_test} verifica el gradiente de constantes.

\subsection{Línea base RK4/FDM}
El mismo solver numérico a baja resolución (50 nodos por eje) se mide contra la referencia. Es la
columna RK4/FDM de las tablas: devuelve una grilla, no una fórmula. Solo para problemas estacionarios.

\subsection{Sondeo del operador (priors)}
\code{probe\_priors} (\code{src/pde\_problems.cpp}) evalúa el residuo sobre \emph{funciones genéricas}
(jets sintéticos con derivadas hasta orden 4) para descubrir propiedades del operador. Usa los datos de
borde solo sobre la frontera; nunca la solución en el interior.

\nota{\textbf{Principio:} los priors \emph{sesgan} la inicialización y la elección de esqueletos, pero
nunca \emph{excluyen} candidatos. Que el operador tenga una propiedad no implica que la solución la
tenga (depende del borde, del forzamiento y de posibles rupturas de simetría).}

{\small
\begin{longtable}{@{}p{4.2cm}p{11cm}@{}}
\toprule\textbf{Prior} & \textbf{Cómo se detecta / para qué se usa}\\\midrule\endhead
Autonomía en $x$, $y$, $t$ & El residuo no cambia al trasladar el punto. Activa esqueletos espectrales.\\
Paridad par en $x$ & Residuo invariante ante $x\to-x$.\\
Polo en el origen & Residuo enorme o no finito en el origen. Activa esqueletos racionales y de
estiramiento logarítmico $\log(x+\varepsilon)$ (Thomas--Fermi, Lane--Emden).\\
Modos exponenciales & Operador lineal, homogéneo y autónomo. Siembra $u=\sum_k A_k e^{\lambda\cdot\mathbf x}$;
con probabilidad $3/4$ una componente de $\lambda$ se despeja de la relación de dispersión
$P(\lambda)=0$ (símbolo del operador, por secante compleja), así que la semilla ya resuelve la EDP y
solo falta ajustar el borde. Esos individuos nacen como U directo.\\
Invarianza de escala & Homogeneidad aproximada. Activa el esqueleto de Frobenius $x^{1/2}\sum x^n$.\\
Orden de derivación & Se anulan los términos de cada orden y se ve si el residuo cambia (lo verifica el oráculo).\\
Simetría especular $x\leftrightarrow1-x$ & Operador invariante ante la reflexión \emph{y} datos de borde
simétricos (evaluados en la frontera). Siembra $N(x(1-x))$.\\
Separabilidad (2D) & Derivada cruzada del residuo $\approx0$ para $\sin x+\sin y$ o $\sin x\sin y$. Siembra $f(x)+g(y)$ o $f(x)g(y)$.\\
Suma modal & Operador afín-lineal (superposición). Siembra $f(v)+g(x)h(y)$ (cubre Kovasznay).\\
Invarianza rotacional & El operador trata $x$ e $y$ igual. Esqueleto de onda plana $f(ax+by)$.\\
Oscilador amortiguado & Término lineal en la primera derivada. Esqueleto $e^{-av}\sin(bv)$.\\
Separabilidad triple & 2D no estacionario: $f(x)g(y)h(t)$.\\\bottomrule
\end{longtable}}

\subsection{Detector de singularidad}
En el constructor de \code{PISolver} se compara el residuo de un árbol de prueba $x+1$ en el centro y
cerca del borde ($10^{-5}$). Si el borde es 1000 veces peor, el problema se marca singular
(\code{Config::IS\_SINGULAR}) y el descenso por gradiente de constantes recorta el gradiente a $\pm50$.

\subsection{Puntos de entrenamiento y de validación}
\begin{itemize}
\item \textbf{Entrenamiento} (al azar, por semilla): 500 puntos de dominio (en 1D con distribución de
  Chebyshev, que concentra puntos en los bordes) y 200 de frontera, muestreados del mismo conjunto que
  describe la validación (incluida la condición inicial si corresponde).
\item \textbf{Validación} (fija y determinista, distinta del entrenamiento): 401 puntos interiores en 1D;
  $41\times21$ en 1D no estacionario; $41\times41$ en 2D; $21\times21\times11$ en 2D no estacionario. La
  frontera de validación cubre aristas (41 o 21 por arista, por instante) y la condición inicial.
\end{itemize}

%=====================================================================
\section{Representación de las soluciones}
%=====================================================================

\subsection{Árboles de expresión}
Cada individuo es un árbol (\code{src/tree\_node.cpp}) con el vocabulario:
\begin{itemize}
\item \textbf{Binarios:} $+,\ -,\ \times,\ \div,\ $ potencia.
\item \textbf{Unarios:} $\sin,\cos,\sinh,\cosh,\exp,\tanh,\log$, cuadrado, gaussiana $e^{-a^2/2}$.
\item \textbf{Polinomios ortogonales:} Legendre, Hermite, Chebyshev, Laguerre (grado $0$--$3$), evaluados
  sobre argumento complejo.
\item \textbf{Series:} $\sum_{n=1}^{5} C_n\,f_n$ (esqueletos espectral y de Frobenius).
\item \textbf{Hojas:} variables $x,y,t$; constantes evolucionables (ERC); $\pi$, $e$, la unidad
  imaginaria $i$; constantes físicas $G, c, \hbar, k_B, \varepsilon_0$ (para el análisis dimensional).
\end{itemize}
Profundidad máxima 6. Los árboles se simplifican algebraicamente después de cada operador genético.

\subsection{Diferenciación automática (jets)}
Cada nodo devuelve un \emph{jet}: el valor y todas las derivadas parciales necesarias, propagadas con las
reglas de Leibniz (producto) y de Faà di Bruno (composición). El jet estándar llega a orden 2
(\code{ad\_eval\_t}); para Navier--Stokes se usa el extendido de orden 4 con derivadas mixtas
(\code{ad\_eval\_ext\_t}). Todo en aritmética compleja, así que Schrödinger y las rotaciones de fase
$e^{i\theta}$ funcionan de forma nativa.

\paragraph{Proyección real.} En problemas de campo real, la solución es $\mathrm{Re}(U)$ y se proyecta el
jet completo \emph{antes} de aplicar el operador (\code{project\_real}): las derivadas conmutan con la
parte real, pero $\mathrm{Re}(U)^3\ne\mathrm{Re}(U^3)$.

\subsection{Ansatz de frontera exacta}
Cuando el problema es estacionario y de orden $\le2$ (\code{ansatz\_applies}), el árbol puede representar
$N$ y la solución se construye como $U = L + B\,N$, con $L$ que interpola los datos y $B$ que se anula
donde hay dato (\code{apply\_boundary\_ansatz}):
\begin{align*}
\text{1D Dirichlet:}\quad & U = g_0 + x(g_1-g_0) + x(x-1)\,N\\
\text{1D valor inicial:}\quad & U = u_0 + u_0'\,x + x^2 N\\
\text{2D Dirichlet:}\quad & U = \mathcal C[g](x,y) + x(1-x)\,y(1-y)\,N
\end{align*}
donde $\mathcal C[g]$ es la interpolación transfinita (de Coons) de los datos en las cuatro aristas. Con
ansatz, el error de frontera es cero por construcción. En los problemas no estacionarios y en
Navier--Stokes el ansatz no aplica, y la frontera es un objetivo medido.

\paragraph{Verificación del ansatz} (\code{ansatz\_border\_check}): la garantía solo vale si $N$ no
diverge cerca del borde (un $N\sim C/x$ deja $BN\to C\ne0$). Se evalúa $U$ a $10^{-6}$ del borde y se
compara con el dato; un desvío $>0.05$ hace al individuo infactible, y el desvío entra en el puntaje de
validación.

\subsection{Representación híbrida}
Con \code{HYBRID\_REPR} activo (por defecto), conviven en la población individuos con ansatz
($\code{raw}=\text{false}$, el árbol es $N$) e individuos directos ($\code{raw}=\text{true}$, el árbol es
$U$ y la frontera se mide). Motivo: una forma cerrada simple como $\sin\pi x$ es corta como $U$ pero
enorme como $N=(U-L)/B$. Un tercio de la población inicial no sesgada nace directa, y cada hijo cambia
de representación con probabilidad 5\%. Puede fallar en algunos problemas (Troesch); se compara con la
ablación \code{-{}-no-hybrid}.

\subsection{La frontera como límite lateral}
Para individuos directos, el dato de borde se exige también como límite desde el interior
(\code{boundary\_error}). Motivo, la \emph{trampa del pico}: un término $e^{-(x/\varepsilon)^2}$ con
$\varepsilon\to0$ fija $U(0)$ exactamente mientras el interior resuelve otra cosa, y ningún punto de
colocación cae dentro del pico. El límite se estima con extrapolación cúbica $3U(d)-3U(2d)+U(3d)$ a dos
escalas, $d=10^{-3}$ y $d=10^{-5}$, y se toma el peor de los errores. En operadores de orden 4 se exige
además la derivada normal $\partial_n U$.

%=====================================================================
\section{Evaluación de un individuo}
%=====================================================================
\code{PIIndividual::evaluate} (\code{src/pi\_solver.cpp}).

\subsection{Factibilidad}
Solo validez de la expresión y gramática del lenguaje; ningún prior es una restricción. Cada violación
suma a \code{constraint\_violation}:
\begin{center}\small
\begin{tabular}{@{}lr@{}}\toprule
Restricción & Peso\\\midrule
No usa todas las variables independientes ($x$; $y$ en 2D; $t$ si no es estacionario) & 10\\
Residuo no finito & 10\\
Dimensión física inconsistente (solo árboles sin ERC) & 5\\
Profundidad unaria $>3$; profundidad $>6$ & 1 c/u\\
Trigonométrica anidada; argumento trigonométrico no afín & 1 c/u\\
Polinomio anidado; exponencial anidada; grado de polinomio inválido & 1 c/u\\
El ansatz no alcanza el dato de borde ($N$ diverge en el borde) & 1\\
La solución no está definida en la frontera (NaN o infinito) & 1\\\bottomrule
\end{tabular}\end{center}
La comparación entre individuos usa la dominancia con restricciones de Deb: un factible domina a un
infactible; entre infactibles gana el de menor violación.

\subsection{Objetivos}
\begin{itemize}
\item $f_1$ (\code{mse\_domain}): promedio ponderado de $|\mathcal N[U]|^2$ en el lote de dominio, con
  recorte por punto en $10^4$ para que una singularidad no domine el promedio. Los pesos por punto se
  explican en la sección~\ref{sec:pesos}.
\item $f_2$ (\code{mse\_boundary}): error de borde medido (cero con ansatz).
\item $f_3$ (\code{tree\_size}): número de nodos.
\end{itemize}
Si todo individuo tiene $f_2=0$ (ansatz sin híbrido), $f_2$ se excluye de NSGA-II: un objetivo empatado
no aporta a la dominancia pero distorsiona la distancia de hacinamiento.

%=====================================================================
\section{Población inicial}
%=====================================================================
\code{PISolver::run}, con tamaño de población 500:
\begin{enumerate}
\item \textbf{Un tercio sesgado por los priors activos}, repartido entre ellos: separable, suma modal,
  simétrico especular, modos exponenciales.
\item \textbf{La mitad del resto} con \code{random\_tree\_special}: esqueletos físicos típicos
  ($1/(1+x^2+y^2)$, gaussiana, $\sin\pi x\sin\pi y$, $\cos\cos$, $\tanh\tanh$, $\sin(\pi x)e^{-t}$,
  $xe^{-x^2}$, \dots) más los activados por los priors: racional, espectral, Frobenius, estiramiento
  logarítmico, onda plana rotada, oscilador amortiguado.
\item \textbf{El resto} con árboles aleatorios: hojas 60\% variables y 40\% constantes; nodos internos
  40\% aritméticos, 40\% unarios y 20\% polinomios ortogonales.
\end{enumerate}
Cada árbol nace multiplicado por una constante evolucionable. Un tercio de los no sesgados nace como
$U$ directo (híbrido).

%=====================================================================
\section{El bucle evolutivo}
%=====================================================================
Hasta 1000 generaciones. Orden de las operaciones en cada generación $g$:

\subsection{Lote de entrenamiento}
\begin{itemize}
\item Mini-lote del 80\% de los puntos disponibles, re-muestreado cada 8 generaciones (no en cada una:
  con un blanco estable la selección premia mejoras reales, no puntos más fáciles).
\item Un lote grande hace improbable que una función degenerada tenga residuo bajo por casualidad
  (se observó con un candidato casi trivial en Airy 1D).
\item En las últimas 50 generaciones se usa la grilla completa.
\end{itemize}

\subsection{Refinamiento adaptativo de residuos por comité (RAR)}
Cada 25 generaciones (\code{apply\_committee\_rar}): un comité formado por las 4 mejores élites y un 40\%
de individuos al azar evalúa 10\,000 candidatos (1D) o 2\,000 (2D). El 30\% de los puntos del conjunto
de entrenamiento se reemplaza por los candidatos con mayor residuo de consenso. Así los puntos difíciles
entran al lote.

\subsection{Pesos por punto}\label{sec:pesos}
Inspirados en el análisis NTK de PINNs (Wang et al., 2022): dentro del lote, cada punto pesa
$1+2\,r_i/\max r$, donde $r_i$ es el residuo promedio de hasta 6 individuos del frente. Rango $[1,3]$:
los puntos fáciles nunca pierden señal y los difíciles pesan hasta el triple. RAR decide \emph{qué}
puntos entran; los pesos, \emph{cuánto} pesa cada uno.

\subsection{Orden no dominado y campeón}
Orden no dominado rápido de NSGA-II con restricciones, y actualización del Hall of Fame
(sección~\ref{sec:campeon}).

\subsection{Parada temprana}
Si el error de entrenamiento del campeón baja de $10^{-6}$, se confirma en la grilla fija de validación
(debe bajar de $10^{-5}$). Si se confirma, la corrida termina.

\subsection{Reproducción}
Selección por torneo de tamaño 4 (gana el de mejor rango y, a igual rango, el de mayor hacinamiento).
Cada hijo:
\begin{itemize}
\item \textbf{Cruzamiento} (prob.\ 0.8): se intercambian términos aditivos entre los padres; si ambos
  términos son productos, con probabilidad 0.5 se intercambian factores (no se destruye la estructura de
  un término que funcionaba).
\item \textbf{Mutación} (\code{tree\_mutate}): 55\% mutación puntual (cambiar el tipo de nodo dentro de su
  familia funcional); si no, sobre los términos aditivos: borrar uno, agregar uno o reemplazar uno.
  El reemplazo regenera el término entero con probabilidad $0.6\times$\emph{agresividad} y, si no, muta
  a nivel de factor (borrar, agregar o reemplazar un factor, o agregar una rotación de fase
  $e^{i\theta}$ con $\theta$ evolucionable).
\item \textbf{Enfriamiento:} agresividad $=0.3+0.7(1-g/G)$. Exploración fuerte al inicio, refinamiento al
  final, con un piso de 0.3 para conservar siempre una vía de escape de óptimos locales.
\item \textbf{Élites} (algún padre en el frente 1): 70\% solo mutación de constantes ($\sigma=0.1$),
  30\% mutación estructural. \textbf{No élites:} mutación estructural con prob.\ 0.45 y de constantes
  ($\sigma=0.2$) con prob.\ 0.4.
\item Cambio de representación ($N\leftrightarrow U$) con prob.\ 0.05, y simplificación algebraica.
\end{itemize}

\subsection{Selección de sobrevivientes}
Padres e hijos se combinan (1000 individuos) y se eligen 500 por rango y distancia de hacinamiento
(\code{nsga2\_select\_next}). Antes se \textbf{eliminan duplicados}: individuos con el mismo tamaño y los
mismos objetivos (a $10^{-12}$) pasan al final y solo entran si faltan únicos. Sin esto, un atractor
trivial (por ejemplo $N=0$, tamaño 1) se clona hasta llenar el frente y congela la búsqueda.

\subsection{Archivo, pulido y registro}
\begin{itemize}
\item \textbf{Archivo histórico:} los 4 números (residuo, frontera, tamaño, rango) de cada individuo de
  cada generación (\code{*\_population\_archive.csv}). En la última generación se guarda además una
  copia sin pulir de la población.
\item \textbf{Pulido de constantes de toda la población:} descenso por gradiente con paso adaptativo
  (lr inicial 0.05; $\times1.1$ si mejora, $\times0.5$ y se revierte si empeora), gradiente por
  diferencias centradas con paso relativo $10^{-6}$ (válido para árboles complejos). 30 iteraciones cada
  10 generaciones y 10 en las demás. Así una estructura buena con constantes mal iniciadas no pierde
  contra una peor con mejor suerte.
\item \textbf{Registro de convergencia} (\code{*\_convergence.csv}): mínimos del frente en el lote,
  puntaje de validación del campeón y error del campeón contra la referencia (solo registro).
\end{itemize}

\nota{Las columnas \code{best\_mse\_domain} y \code{best\_mse\_boundary} son mínimos \emph{por separado}
sobre el frente, de individuos distintos. En ecuaciones donde una constante es solución (Airy, Fisher,
Ginzburg--Landau, Lane--Emden, Troesch, NS, Schrödinger no estacionario), $u\equiv0$ vive en el frente
como extremo de residuo cero, así que esa columna vale 0. Para graficar se usan
\code{best\_val\_score} y \code{champion\_ref\_err}.}

%=====================================================================
\section{Selección y pulido del campeón}\label{sec:campeon}
%=====================================================================

\paragraph{Hall of Fame} (\code{update\_hall\_of\_fame}). En cada generación:
\begin{enumerate}
\item Entre los individuos del frente 1 factibles y que usan todas las variables, se toma el de menor
  error de entrenamiento (ya calculado, sin costo extra).
\item Se compara con el campeón vigente mediante el \textbf{puntaje de validación}: residuo y frontera en
  las grillas fijas, más el desvío del ansatz en el borde. Solo información de la EDP; la referencia
  nunca entra.
\end{enumerate}

\paragraph{La solución trivial.} $u\equiv0$ puede estar en el frente (residuo cero), pero su puntaje de
validación incluye el error de frontera, así que nunca gana. Es el modo de falla que Rockett (2026) deja
abierto.

\paragraph{Pulido final.} Al terminar, el campeón se pule con 200 iteraciones de descenso por gradiente
contra la grilla fija de validación (no contra el mini-lote, para no sobreajustar). Luego se mide su
error contra la referencia (\code{solution\_mse}), que es la métrica de las tablas.

%=====================================================================
\section{Salidas y post-proceso}
%=====================================================================

\subsection{Archivos por problema y semilla}
\begin{center}\small
\begin{tabular}{@{}p{6cm}p{9cm}@{}}\toprule
Archivo & Contenido\\\midrule
\code{<P>\_pi\_gn\_pareto.csv} & Población final (todos los rangos).\\
\code{<P>\_pi\_gn\_pareto\_unpolished.csv} & Población final antes del pulido.\\
\code{<P>\_pi\_gn\_population\_archive.csv} & Archivo histórico ($\sim$340 MB por semilla; no va a git).\\
\code{<P>\_pi\_gn\_convergence.csv} & Registro por generación.\\
\code{expr\_<P>\_PISR-EMOAD.tex} / \code{.sym} & Fórmula del campeón (LaTeX y SymPy, con \code{raw}, \code{uses\_ansatz}).\\
\code{grid\_<P>\_PISR-EMOAD.csv} & Campeón y referencia en una grilla (t = 0.5 si no es estacionario).\\
\code{comparison\_summary.csv} & Resumen por EDP de la semilla.\\
\code{results/all\_runs\_summary.csv} & Todas las semillas (fusión por EDP; \code{-{}-only} reemplaza solo su fila).\\\bottomrule
\end{tabular}\end{center}

\subsection{Reporte} (\code{./generate\_report.sh})
\begin{enumerate}
\item \code{consolidate\_runs.py}: arma \code{all\_runs\_summary.csv} (conserva el hipervolumen).
\item Figuras: soluciones, frentes, convergencia (mediana y banda intercuartil sobre las semillas).
\item \code{simplify\_expressions.py}: simplifica la fórmula con SymPy, redondea constantes a valores
  exactos ($p/q$, $\pi$, $\sqrt2$, \dots) y verifica contra la grilla del campeón (tolerancia relativa
  $\le10^{-3}$).
\item \code{extrapolation.py}: evalúa el campeón en $[0,1.5]$ y lo compara con un polinomio de Chebyshev
  de igual precisión interior.
\item \code{physics\_checks.py}: extrae cantidades físicas de la fórmula (energía, $u'(0)$, número de
  onda, dispersión, radio de Lane--Emden, tasa de decaimiento).
\item \code{interpretability.py}: PISR contra PINN+PySR (error, nodos, recuperación exacta,
  extrapolación).
\item \code{generate\_report.py} y \code{pdflatex}.
\end{enumerate}
\nota{\textbf{Semilla mediana:} todo lo que necesita una sola fórmula (expresiones, gráficos por
problema, extrapolación, física) usa la semilla cuyo error es el 15.º de 30. Las tablas de error
agregan todas las semillas.}

\subsection{Competidores} (\code{./run\_competitors.sh -{}-runs 30})
\begin{itemize}
\item \code{PISR-NSGA-II -{}-export-reference} escribe, por EDP, los puntos de validación, la referencia y
  los datos de borde (\code{results/reference/}). Así los competidores se miden en los mismos puntos y
  con la misma referencia.
\item \textbf{PINN} (\code{competitors/pinn.py}): red $4\times50$ tanh en float64, con los mismos residuos
  (copiados de \code{pde\_residual\_ad} y verificados contra las soluciones exactas). 10\,000
  iteraciones de Adam y 2\,000 de L-BFGS, con tope de 30 min por trabajo.
\item \textbf{PINN+PySR} (\code{competitors/pinn\_pysr.py}): PySR ajusta una fórmula a la predicción del
  PINN (nunca a la referencia). Guarda la fórmula elegida y el frente completo.
\end{itemize}

%=====================================================================
\section{Dónde se usa la referencia (garantías contra fugas)}
%=====================================================================
\begin{center}\small
\begin{tabular}{@{}p{6.5cm}cp{6.5cm}@{}}\toprule
Componente & ¿Usa la referencia? & Qué usa\\\midrule
Evaluación, objetivos, RAR, pesos & No & Residuo y datos de borde\\
Sondeo de priors & No & Operador sobre funciones genéricas; borde solo en la frontera\\
Hall of Fame, parada temprana, pulido & No & Puntaje de validación (física)\\
Registro de convergencia & Solo escritura & \code{champion\_ref\_err}, que nada lee\\
\code{solution\_mse}, reporte, oráculo & Sí & Evaluación después de la búsqueda\\\bottomrule
\end{tabular}\end{center}

%=====================================================================
\section{Parámetros}
%=====================================================================
\begin{center}\small
\begin{tabular}{@{}lll@{}}\toprule
Parámetro & Valor & Opción de línea de comandos\\\midrule
Población / generaciones & 500 / 1000 & \code{-{}-pop}, \code{-{}-gen}\\
Puntos de dominio / frontera & 500 / 200 & \code{-{}-domain}, \code{-{}-boundary}\\
Profundidad máxima & 6 & \code{-{}-depth}\\
$\sigma$ de mutación de constantes & 0.20 & \code{-{}-sigma}\\
Prob.\ de cruzamiento / mutación & 0.80 / 0.45 & ---\\
Torneo & 4 & ---\\
Umbral de parada & $10^{-6}$ (validación $10^{-5}$) & \code{-{}-stop}\\
RAR: intervalo / candidatos / fracción & 25 / 10\,000 / 0.30 & ---\\
RAR: élites / fracción al azar & 4 / 0.40 & ---\\
Mini-lote / re-muestreo & 80\% / cada 8 gen. & ---\\
Representación híbrida & activa & \code{-{}-no-hybrid}\\
Ansatz de frontera & activo & \code{-{}-no-ansatz}\\\bottomrule
\end{tabular}\end{center}

%=====================================================================
\section{Resultados de las 30 semillas (resumen)}
%=====================================================================
\begin{itemize}
\item 15 de 21 problemas bajo $10^{-3}$ en al menos el 83\% de las semillas.
\item Recuperación exacta de la forma cerrada: Laplace, Poisson, oscilador y Schrödinger 1D (30/30);
  Schrödinger 2D (29/30); Schrödinger no estacionario 1D (26/30); NS no estacionario (12/30);
  Ginzburg--Landau (11/30; las demás son aproximaciones de otra forma con error $\sim10^{-8}$).
\item No resueltos: Thomas--Fermi 1D y 2D (singularidad en el origen), Troesch (capa límite rígida),
  NS Kovasznay (no linealidad de orden 4), Schrödinger no estacionario 2D (la búsqueda mejora el
  residuo sin acercarse a la solución).
\item Tiempo mediano: unas 3.2 h por semilla para la suite completa en 32 núcleos.
\end{itemize}

%=====================================================================
\section{Historial de correcciones importantes}
%=====================================================================
Para recordar por qué el código es como es:
\begin{itemize}
\item \textbf{Problemas redefinidos} (20 de octubre de 2026): Airy verdadero, Lane--Emden como IVP,
  Thomas--Fermi en $[0,1]$, Schrödinger no resonante, Ginzburg--Landau con $\tanh(x/\sqrt2)$, NS con el
  residuo real. Los resultados anteriores a esa fecha se excluyen del reporte.
\item \textbf{La referencia salió de la selección:} el Hall of Fame y la parada temprana comparaban
  contra la referencia; ahora usan solo el puntaje físico de validación.
\item \textbf{Trampa del pico:} cerrada exigiendo la frontera como límite lateral.
\item \textbf{Laguerre con argumento complejo:} los polinomios se evaluaban sobre la parte real, y
  Lane--Emden daba residuo 0 con $u\equiv1$.
\item \textbf{Impresión de fórmulas:} $A-(B+C)$ salía sin paréntesis y las constantes físicas salían
  vacías; se corrigió y se recuperaron los signos de las fórmulas antiguas.
\item \textbf{Kovasznay recalibrado} a $\mathrm{Re}=20$ con amplitud $1/(2\pi)$; con $\mathrm{Re}=100$ el
  problema era casi trivial ($\psi\approx y$ ya daba $10^{-7}$).
\item \textbf{Resúmenes que se pisaban:} \code{-{}-only} borraba las demás filas; ahora se fusionan por EDP.
\end{itemize}

\end{document}
